%% ma: L10-B01-0035
%% lop: 10
%% chuong: 1
%% bai: 1
%% dang: 10.1.5
%% loai: TLN
%% muc: TH
%% dap_an: "2"
%% nguon: "Ảnh giáo viên gửi 10/10/2026 (ThS. Nguyễn Hoàng Việt, Chương 1 Mệnh đề), B-Đề 2, Phần 3 Câu 19"
%% kien_thuc: "Bài 1 (mệnh đề phủ định)"
%% trang_thai: cho-duyet
%% ly_do: "Dạng toán mới đề xuất 10.1.5: mệnh đề phủ định"
%% ghi_chu: "Đề gốc không có đáp án. Đáp số 2 theo định nghĩa phân số có tử và mẫu nguyên (ý a sai); nếu coi 5/1,2 là phân số thì đáp số là 3. Giáo viên kiểm tra lại"
\begin{ex}
Lập mệnh đề phủ định của mỗi mệnh đề sau và cho biết có bao nhiêu mệnh đề phủ định sai.
\begin{enumerate}[label=\alph*)]
\item $A\colon$ ``$\dfrac{5}{1{,}2}$ là một phân số''.
\item $B\colon$ ``Phương trình $x^2+3x+2=0$ có nghiệm''.
\item $C\colon$ ``$2^2+2^3=2^{2+3}$''.
\item $D\colon$ ``Số $2025$ chia hết cho $15$''.
\end{enumerate}
\shortans{2}
\loigiai{Mệnh đề phủ định sai khi mệnh đề gốc đúng. $A$ sai (phân số có tử và mẫu là số nguyên, mẫu $1{,}2$ không nguyên) nên $\overline A$ đúng. $B$ đúng ($x=-1$) nên $\overline B$ sai. $C$ sai ($12\ne32$) nên $\overline C$ đúng. $D$ đúng ($2025=15\cdot135$) nên $\overline D$ sai. Có $2$ mệnh đề phủ định sai.}
\end{ex}
